\begin{textsample}{POS Dim 4 – 2021\_c5 – Score 1.00 – 2021\_c5/enade\_2021\_tec\_gest\_tec\_inf\_2a1ac18999.txt}  \label{ex:f4_pos_213}
CAPÍTULO 5 DISTRIBUIÇÃO DOS CONCEITOS 5.1 PANORAMA NACIONAL DA DISTRIBUIÇÃO DOS CONCEITOS O Conceito Enade é um indicador de qualidade que avalia os cursos por intermédio dos desempenhos dos estudantes a partir dos resultados do Enade.
Os valores possíveis vão de 1 ( pior situação ) a 5 ( melhor situação ).
Os cursos com apenas um ou sem nenhum concluinte participante não obtêm o Conceito Enade, ficando Sem Conceito ( SC ).
Na Tabela 5.1, são apresentadas a quantidade e a distribuição de cursos de Tecnologia em Gestão da Tecnologia da Informação participantes do Enade/2021, por faixa de conceito e Grande Região.
Enfatiza -se, mais uma vez, que a diferença entre os cursos tabulados neste capítulo e no capítulo 2 corresponde aos cursos que não foram avaliados, ou seja, àqueles, em princípio, sem alunos concluintes que tenham sido inscritos para a prova.
Observando -se os dados da Tabela 5.1, nota -se que, dos 148 cursos participantes, 55 ( 37,2 \% ) classificaram -se com conceito 3, o valor modal.
Esse foi também o conceito modal na maioria das regiões : Nordeste ( 47,4 \% ), Sudeste ( 36,5 \% ), Sul ( 42,9 \% ) e Norte ( 33,3 \%, em uma distribuição multimodal, com a mesma proporção sendo atribuída aos conceitos 2 e 4 ).
O conceito 2 foi o segundo mais frequente em nível nacional ( 29,7 \%, correspondendo a 44 cursos ), tendo sido o conceito modal da região Centro-Oeste ( 46,2 \% ), e o conceito 4, o terceiro ( 17,6 \%, correspondendo a 26 cursos ).
Houve, ainda, 13 cursos ( 8,8 \% ) que receberam o conceito 1, e nove cursos ( 6,1 \% ) que receberam o conceito 5.
Dos 148 cursos de Tecnologia em Gestão da Tecnologia da Informação, um ( 0,7 \% ) ficou Sem Conceito ( SC ).
A região Norte participou com três cursos, o correspondente a 2,0 \% do total nacional : os conceitos foram o 2, o 3 e o 4 ( um curso em cada conceito ).
A região Nordeste participou com 19 cursos, correspondendo a 12,8 \% do total nacional.
Nessa região, nove cursos ( 47,4 \% em termos regionais ) obtiveram conceito 3, o conceito modal para a região, como já comentado.
Os conceitos 1, 4 e 5 foram atribuídos a dois cursos cada, e quatro cursos receberam o conceito 2.
Nenhum curso ficou Sem Conceito ( SC ).
Dos 85 cursos participantes da região Sudeste ( 57,4 \% do total nacional ), o conceito 3 foi o conceito modal, como já anteriormente assinalado, tendo sido obtido por 31 cursos ( 36,5 \% ).
O conceito 1 foi atribuído a nove cursos ( 10,6 \% ), e o conceito 2, a 28 cursos ( 32,9 \% ).
Receberam os conceitos 4 e 5, respectivamente, 14 cursos ( 16,5 \% ) e dois cursos ( 2,4 \% ).
Nessa região, um curso ( 1,2 \% ) ficou Sem Conceito ( SC ).
Os 28 cursos da região Sul corresponderam a 18,9 \% do total nacional.
Como já indicado, houve predominância do conceito 3, atribuído a 42,9 \% dos cursos da região ( conceito modal ), ou seja, atribuído a 12 dos 28 cursos participantes na região Sul.
O conceito 1 foi atribuído a um curso ( 3,6 \% ) ; o conceito 2, a cinco cursos ( 17,9 \% ) ; e os conceitos 4 e 5 foram atribuídos a seis ( 21,4 \% ) e a quatro cursos ( 14,3 \% ), respectivamente.
Nenhum curso ficou Sem Conceito ( SC ) na região Sul.
Dos 13 cursos participantes na região Centro-Oeste ( 8,8 \% do total nacional ), seis ( 46,2 \% em termos regionais ) receberam conceito 2, o conceito modal, como já destacado.
Um curso ( 7,7 \% ) obteve o conceito 1 ; um outro curso, o conceito 5 ; dois cursos ( 15,4 \% ) obtiveram o conceito 3 ; e três cursos ( 23,1 \% ) ; o conceito 4.
Nessa região, nenhum curso ficou Sem Conceito ( SC ).
No Gráfico 5.1, é apresentada a distribuição cumulativa do Conceito Enade, segundo a Grande Região.
Em um gráfico de distribuição cumulativa, em cada valor de abscissa ( eixo x ), o valor apresentado no eixo das ordenadas ( eixo y ) é o valor acumulado da variável ( Conceito Enade ).
Observando -se a linha vertical pontilhada, que ajuda a visualizar o percentual de conceitos acumulados até o valor 3, ou seja, os conceitos 1, 2 e 3, pode -se notar que, por exemplo, a região Sudeste ( linha vermelha ) apresentou 81,0 \% dos cursos com o conceito 3 ou abaixo deste valor.
A região Sul ( linha rosa ), por outro lado, apresentou 64,3 \% dos cursos na mesma situação : com o conceito 3 ou abaixo.
Destaca -se que, em algumas regiões, o conceito 5 pode não ter sido alcançado por nenhum curso.
Foi o que ocorreu na região Norte ( linha azul ), na qual todos os cursos obtiveram conceitos menores ou iguais a 4, tendo a poligonal alcançado 100 \% no conceito 4.
Em linhas gerais, em um gráfico de distribuição cumulativa, poligonais mais à esquerda representam as Grandes Regiões com a pior distribuição, e poligonais mais à direita, as Grandes Regiões com os melhores conceitos.
A região Sul ( linha rosa ) apresentou a Cursos em IES situadas nas regiões Sudeste ( linha vermelha ) e Centro-Oeste ( linha laranja ) apresentaram uma situação pior ( linhas mais à esquerda ), pelo menos em parte dos conceitos, mas fora a região Sul, as demais poligonais se entrecortam e não se pode ordenar as distribuições. 5.2 CONCEITOS POR CATEGORIA ADMINISTRATIVA, POR MODALIDADE DE ENSINO E GRANDE REGIÃO Na Tabela 5.2, é apresentada a distribuição dos cursos de Tecnologia em Gestão da Tecnologia da Informação participantes do Enade/2021, por Categoria Administrativa e Modalidade de Ensino, de acordo com os conceitos por eles alcançados, segundo as Grandes Regiões brasileiras.
Dos 148 cursos participantes, 19 ( 12,8 \% ) eram ministrados em Instituições Públicas, e 129 ( 87,2 \% ), em Instituições Privadas.
Quanto à Modalidade de Ensino, 102 cursos eram Presenciais, e 46, a Distância.
De acordo com as informações da Tabela 5.2, em termos nacionais, entre Instituições Públicas e Privadas, observa -se que, dos nove cursos avaliados com o conceito 5, dois foram oferecidos em IES Públicas, e sete, em IES Privadas.
Dos 19 cursos participantes de IES Públicas, o conceito 4 foi o valor modal, atribuído a oito cursos ( 42,1 \% ).
Dos demais cursos de Instituições Públicas participantes, quatro cursos ( 21,1 \% ) foram avaliados com o conceito 2 ; cinco cursos ( 26,3 \% ), com o conceito 3, e, como já comentado, dois cursos ( 10,5 \% ), com o conceito 5.
Nessa categoria, nenhum curso ficou Sem Conceito ( SC ) ou recebeu o conceito 1.
Na rede privada, o conceito modal foi o 3, atribuído a 50 cursos, o correspondente a 38,8 \% dos 129 cursos da categoria.
Dos demais cursos participantes, 13 ( 10,1 \% ) receberam o conceito 1, e 40 ( 31,0 \% ), o conceito 2.
O conceito 4 foi atribuído a 18 cursos ( 14,0 \% ), e o conceito 5, a sete cursos ( 5,4 \% ).
Nessa Categoria Administrativa, um curso ( 0,8 \% ) ficou Sem Conceito ( SC ).
Para os 102 cursos participantes da Modalidade Presencial, o conceito 2 foi o valor modal, atribuído a 37 cursos ( 36,3 \% ).
Dos demais cursos nessa Modalidade de Ensino, 13 cursos ( 12,7 \% ) receberam o conceito 1 ; 33 cursos ( 32,4 \% ) receberam o conceito 3 ; 12 cursos ( 11,8 \% ), o conceito 4 ; e seis cursos ( 5,9 \% ), o conceito 5.
Nessa Modalidade de Ensino, um curso ( 1,0 \% ) ficou Sem Conceito ( SC ).
Foram 46 os cursos a Distância informados, dos quais 22 ( 47,8 \% ) receberam o conceito modal 3.
O conceito 2 foi atribuído a sete cursos ( 15,2 \% ) ; o conceito 4, a 14 cursos ( 30,4 \% ) ; e o conceito 5, a três cursos ( 6,5 \% ).
Nessa Modalidade de Ensino, nenhum curso ficou Sem Conceito ( SC ) ou recebeu o conceito 1.
Na análise por região, observa -se que, na região Norte, as Instituições Privadas participaram com dois cursos ( 66,7 \% do total regional ), que receberam os conceitos 3 e 4.
As Instituições Públicas participaram com um curso na região Norte ( 33,3 \% do total regional ), ao qual foi atribuído o conceito 2.
Todos os cursos oferecidos na região Norte foram Presenciais.
Na região Nordeste, a rede privada concentrou 17 dos 19 cursos participantes, o equivalente a 89,5 \% do total da região.
Entre os cursos oferecidos por IES Privadas no Nordeste, sete cursos foram avaliados com conceito 3, o conceito modal.
Dois cursos obtiveram o conceito 1 ; quatro receberam o conceito 2 ; dois cursos, o conceito 4 ; e dois cursos, o conceito 5.
Nessa combinação de categoria e região, nenhum curso ficou Sem Conceito ( SC ).
As Instituições Públicas da região Nordeste participaram com dois cursos ( 10,5 \% do total da região ), que receberam o conceito 3.
Na região Nordeste, a maioria dos cursos foi oferecida na modalidade Presencial : 15 de 19 cursos, tendo sete cursos recebido o conceito modal 3.
Os demais cursos presenciais da região receberam os conceitos 1 ( dois cursos ), 2 ( quatro cursos ) e 5 ( dois cursos ).
Nenhum curso ficou Sem Conceito ( SC ) ou recebeu o conceito 4 na Modalidade Presencial na região Nordeste.
Dos quatro cursos a Distância participantes da região Nordeste, dois receberam o conceito 3, e dois, o conceito 4.
Na região Sudeste, a proporção de cursos da rede privada, 83,5 \%, correspondeu a 71 dos 85 cursos participantes.
Nessa categoria e região, o conceito modal foi o 3, atribuído a 28 cursos, tendo um curso ficado Sem Conceito ( SC ).
Os demais foram avaliados com os seguintes conceitos : 1 ( nove cursos ), 2 ( 26 cursos ), 4 ( seis cursos ) e 5 ( um curso ).
Para os 14 cursos oferecidos por Instituições Públicas na região Sudeste, o conceito modal foi o 4, atribuído a oito cursos.
Nessa região, dois receberam o conceito 2, três receberam o conceito 3, e um recebeu o conceito 5.
Nenhum curso ficou Sem Conceito ( SC ) ou recebeu o conceito 1.
Também na região Sudeste, a maior parte dos cursos foi oferecida na Modalidade Presencial : 64 de 85 cursos.
Para esses, a moda foi o conceito 2, com 23 cursos.
Os demais cursos Presenciais da região receberam os conceitos 1 ( nove cursos ), 3 ( 21 cursos ), 4 ( nove cursos ) e 5 ( um curso ).
Ficou Sem Conceito ( SC ) um curso.
Na região Sudeste, foram oferecidos 21 cursos na Modalidade Educação a Distância.
O conceito modal foi o 3, atribuído a dez cursos.
Os demais foram avaliados com os conceitos 2 ( cinco cursos ), 4 ( cinco cursos ) e 5 ( um curso ).
Na região Sudeste, nenhum curso da Modalidade a Distância ficou Sem Conceito ( SC ) ou recebeu o conceito 1.
As Instituições Privadas concentraram 27 dos 28 cursos participantes da região Sul, 96,4 \% do total regional.
Desses, 12 cursos obtiveram conceito 3, o conceito modal.
Nessa combinação de Categoria Administrativa e Grande Região, um curso recebeu o conceito 1 ; cinco cursos receberam o conceito 2, outros seis, o conceito 4 ; e ainda três cursos, o conceito 5.
Nenhum curso ficou Sem Conceito ( SC ).
As Instituições Públicas na região Sul participaram com um curso, que recebeu o conceito 5.
Dos 28 cursos participantes da região Sul, nove eram Presenciais.
O conceito modal para essa Modalidade de Ensino foi o 2, atribuído a três cursos.
O conceito 1 foi atribuído a um curso ; o conceito 3, a dois cursos ; o conceito 4, a um curso ; e o conceito 5, a dois cursos.
Nenhum curso Presencial da região Sul ficou Sem Conceito ( SC ).
Já dos 19 cursos oferecidos nessa região na Modalidade a Distância, dez receberam o conceito modal 3.
Nessa Modalidade, o conceito 2 foi atribuído a dois cursos ; o conceito 4, a cinco cursos ; e o conceito 5, a dois cursos.
Nenhum curso ficou Sem Conceito ( SC ) ou recebeu o conceito 1.
Na região Centro Oeste, 12 dos 13 cursos participantes foram de instituições Privadas ( 92,3 \% em termos regionais ).
Desses, cinco cursos receberam o conceito modal 2.
Dos demais cursos, um recebeu o conceito 1 ; um recebeu o conceito 5 ; dois receberam o conceito 3 ; e três, o conceito 4.
O único curso oferecido por Instituições Públicas na região Centro Oeste ( 7,7 \% do total regional ) recebeu o conceito 2.
Na região Centro-Oeste, participaram 11 cursos Presenciais, e o conceito modal foi o 2, recebido por seis cursos.
Quanto aos demais cursos, o conceito 3 foi atribuído a dois cursos, e os conceitos 1, 4 e 5 foram atribuídos a um curso cada.
Nenhum curso ficou Sem Conceito ( SC ).
Os dois cursos a Distância da região Centro-Oeste receberam o conceito 4.
No Gráfico 5.2, apresenta -se a distribuição do Conceito Enade, segundo a Categoria Administrativa da IES.
Os cursos em IES Públicas ( linha azul ) apresentam uma poligonal mais à direita do que os em IES Privadas ( linha verde ) e, consequentemente, uma distribuição de Conceito Enade com valores maiores. 5.3 CONCEITOS POR ORGANIZAÇÃO ACADÊMICA E GRANDE REGIÃO Na Tabela 5.3, encontra -se a distribuição dos conceitos atribuídos aos cursos participantes do Enade/2021, na Área de Tecnologia em Gestão da Tecnologia da Informação, por Organização Acadêmica, segundo as Grandes Regiões brasileiras.
Dos 148 cursos de Tecnologia em Gestão da Tecnologia da Informação participantes, 54 eram oferecidos em Universidades ; 50, em Centros Universitários ; 39, em Faculdades ; e cinco, em CEFET/IF.
Esta distribuição corresponde a, respectivamente, 36,5 \%, 33,8 \%, 26,4 \% e 3,4 \% dos cursos.
De acordo com os dados apresentados, de todos os nove cursos avaliados com o conceito 5, dois eram vinculados a Universidades.
Os cursos vinculados a esse tipo de Organização Acadêmica tiveram o conceito 3 como conceito modal, atribuído a 18 cursos ( 33,3 \% ).
Os demais cursos vinculados a Universidades receberam os conceitos 1 ( oito cursos ), 2 ( 17 cursos ), 4 ( nove cursos ) e 5 ( dois cursos, como já mencionado ).
Nenhum curso ficou Sem Conceito ( SC ).
Dos cursos em Centros Universitários, o conceito modal também foi o 3, atribuído a 26 cursos ( 52,0 \% ).
Nenhum curso vinculado a esse tipo de Organização Acadêmica ficou Sem Conceito ( SC ), e os outros cursos receberam os conceitos 1 ( um curso ), 2 ( 13 cursos ), 4 ( sete cursos ) e 5 ( três cursos ).
Dos 39 cursos mantidos por Faculdades, 12 ( 30,8 \% ) receberam conceito 2, o conceito modal.
O segundo conceito mais frequentemente atribuído aos cursos oferecidos por essa categoria de IES foi o conceito 3, recebido por dez cursos.
Ainda nesse tipo de Organização Acadêmica, quatro cursos receberam o conceito 1 ; nove cursos, o conceito 4 ; e três cursos, o conceito 5.
Dos cursos oferecidos em Faculdades, um ficou Sem Conceito ( SC ).
Foram cinco os cursos mantidos por CEFET/IF, e o conceito modal para esse tipo de Organização Acadêmica foi o 2 ( dois cursos ).
Os conceitos 3, 4 e 5 foram atribuídos a um curso cada.
Nenhum curso ficou Sem Conceito ( SC ) ou recebeu o conceito 1.
Os Centros Universitários da região Norte foram representados por um curso, o qual obteve o conceito 3.
As Faculdades participaram com um curso na região Norte, o qual obteve o conceito 4.
Os CEFET/IF participaram, também, com um curso, o qual obteve o conceito 2.
Na região Norte, não houve cursos mantidos por Universidades.
Na região Nordeste, as Universidades participaram com três dos 19 cursos oferecidos na Área de Tecnologia em Gestão da Tecnologia da Informação.
O conceito modal foi o 4, atribuído a dois cursos.
O terceiro curso recebeu o conceito 2.
Os Centros Universitários contaram com nove cursos participantes na região Nordeste, seis dos quais receberam o conceito modal 3.
Para os demais cursos, foram atribuídos os conceitos 1, 2 e 5 a um curso cada.
Na região Nordeste, nenhum curso mantido por Centros Universitários, ficou Sem Conceito ( SC ) ou recebeu o conceito 4.
As Faculdades foram representadas por seis cursos na região Nordeste, e a distribuição foi bimodal, com os conceitos 2 e 3 sendo atribuídos a dois cursos cada.
Os demais cursos receberam os conceitos 1 e 5 ( um curso cada ).
Na região Nordeste, nenhum curso mantido por Faculdades ficou Sem Conceito ( SC ) ou recebeu o conceito 4.
O único curso oferecido por CEFET/IF na região Nordeste recebeu o conceito 3.
Na região Sudeste, as Universidades concentraram 43 dos 85 cursos de Tecnologia em Gestão da Tecnologia da Informação da região.
Dos cursos oferecidos em Universidades, nessa região, o conceito modal foi o 2, atribuído a 16 cursos.
Os demais cursos receberam os conceitos 1 ( oito cursos ), 3 ( 15 cursos ), 4 ( três cursos ) e 5 ( um curso ).
Nenhum curso ficou Sem Conceito ( SC ).
Os Centros Universitários participaram com 20 cursos na região Sudeste, dos quais dez obtiveram o conceito modal 3.
O conceito 2 foi atribuído a sete cursos, e o conceito 4, a três cursos.
Nenhum curso ficou Sem Conceito ( SC ) nessa combinação de Organização Acadêmica e região ou recebeu os conceitos 1 ou 5.
Já as Faculdades foram representadas por 21 cursos na região Sudeste, dos quais um ficou Sem Conceito ( SC ).
O conceito modal 4 foi atribuído a sete cursos.
Aos demais cursos foram atribuídos os conceitos 1 ( um curso ), 2 ( cinco cursos ), 3 ( seis cursos ) e 5 ( um curso ).
O único curso oferecido por CEFET/IF na região Sudeste recebeu o conceito 4.
Dos 28 cursos da região Sul, sete eram mantidos por Universidades, para os quais a distribuição foi bimodal, com três cursos recebendo o conceito 3, e outros 3 cursos, o conceito 4.
O curso restante recebeu o conceito 5.
Nesse tipo de organização, nenhum curso da região Sul ficou Sem Conceito ( SC ) ou recebeu os conceitos 1 ou 2.
Dos 16 cursos participantes de Centros Universitários da região Sul, nenhum ficou Sem Conceito ( SC ) ou recebeu o conceito 1.
Os conceitos restantes foram atribuídos a três cursos ( conceito 2 ), oito cursos ( conceito modal 3 ), três cursos ( conceito 4 ) e dois cursos ( conceito 5 ).
Quanto aos quatro cursos vinculados a Faculdades na região Sul, o conceito modal foi o 2, atribuído a dois cursos, e os conceitos 1 e 3 foram atribuídos a um curso cada.
O único curso oferecido por CEFET/IF na região Sul recebeu o conceito 5.
Na região Centro-Oeste, apenas um dos 13 cursos oferecidos eram mantidos por Universidades, tendo recebido o conceito 4.
Os Centros Universitários da região Centro-Oeste contaram com quatro cursos, tendo dois cursos sido avaliados com o conceito 2 ( modal ), um curso, com o conceito 3, e um curso, com o conceito 4.
Dos sete cursos oferecidos em Faculdades na região Centro-Oeste, três receberam o conceito 2 ( modal ).
Os conceitos 1, 3, 4 e 5 foram atribuídos a um curso cada.
Nessa região, nenhum curso oferecido em Faculdades ficou Sem Conceito ( SC ).
O único curso oferecido por CEFET/IF na região Centro-Oeste recebeu o conceito 2.
No Gráfico 5.3, apresenta -se a distribuição do Conceito Enade dos cursos de Tecnologia em Gestão da Tecnologia da Informação, segundo a Organização Acadêmica da IES.
Os cursos em CEFET/IF ( linha laranja ) apresentam uma poligonal mais à direita do que as demais, denotando uma melhor distribuição de Conceitos Enade, pelo menos para os maiores valores de conceito.
No outro extremo, encontram -se os cursos mantidos por Universidades ( linha azul ), apresentando a respectiva poligonal mais à esquerda e, consequentemente, uma distribuição de Conceitos Enade com menores valores.
Para valores menores do conceito são os Centros Universitários ( linha verde ) que apresentam a melhor distribuição.
Os cursos mantidos por Faculdades ( linha vermelha ) apresentam poligonais intermediárias e, consequentemente, uma distribuição de Conceitos Enade com valores compatíveis com esta situação.

% matched lemmas: 
\end{textsample}
